\documentclass[10pt,a4paper]{amsart}
\usepackage[margin=19mm]{geometry}
\usepackage{amsmath,amssymb,mathtools}
\usepackage[scr=rsfs,cal=euler]{mathalfa}
\usepackage{xfrac}
\usepackage[unicode,hidelinks,bookmarks=false]{hyperref}
\newcommand{\ie}{i.e.\ }
\newcommand{\cf}{cf.\ }
\newcommand{\resp}{respectively\ }
\newcommand{\Subsection}{\subsection}
\providecommand{\nobreakdash}{\nobreak\hbox{-}}
\setlength{\emergencystretch}{2em}
\multlinegap0pt
\usepackage{tikz}
\usetikzlibrary{cd}
\usepackage{mathtools}
\usepackage{amssymb}
\usepackage[normalem]{ulem}
\usepackage[shortlabels]{enumitem}
\usepackage{xcolor}
\usepackage{tikz}
\usepackage{braket}
\newtheorem{theorem}{Theorem}
\newcommand{\BG}{{\mathbb{G}}}
\newcommand{\BK}{{\mathbb{K}}}
\newcommand{\CA}{{\mathcal A}}
\newcommand{\Ext}{\mathcal{E}\text{xt}}
\newcommand{\FM}{{\mathfrak{M}}}
\newcommand{\rig}{\mathrm{rig}}
\newcommand{\rk}{\mathrm{rk}}
\title{Generalized $K$-theoretic invariants and wall-crossing via non-abelian localization}
\author{Ivan Karpov, Miguel Moreira}
\date{Source: arXiv:2512.22360v2 (CC BY 4.0)}
\begin{document}
\maketitle

\noindent\emph{Source and licence.} Generalized $K$-theoretic invariants and wall-crossing via non-abelian localization. Version arXiv:2512.22360v2 (CC BY 4.0), distributed by arXiv under the Creative Commons Attribution 4.0 International licence, http://creativecommons.org/licenses/by/4.0/, as recorded in the arXiv licence metadata for this exact version and shown on the abstract page https://arxiv.org/abs/2512.22360v2. Full licence notice: \texttt{licenses/arxiv-2512.22360-CC-BY-4.0.txt}.\par\smallskip

\noindent\textit{Editorial scope.} Counted unit: the article's theorem associativity (source0.tex lines 1321--1323). The article's complete original proof of it is delivered with it (source0.tex lines 1324--1357, one \texttt{proof} environment of 34 lines). No mathematics is written, repaired or re-derived: the statement and the proof are the article's own lines, in the article's own notation, and no retained line is substituted or re-flowed.  No further source-local context is needed: every cross-reference of every retained line resolves inside the two delivered spans.  Nothing was omitted from any retained span, so the package carries no omission record.  No retained line contains a citation, so the wrapper carries no bibliography and no bibliography entry.  The retained lines contain the article's own diagram code, which the wrapper typesets with the standard tikz packages; no image file is delivered and the package declares no asset dependency.  Editorial formatting only, in addition: an AMS \texttt{amsart} preamble that restates the article's own declarations and macros used by the retained lines, namely the environments \texttt{theorem} on separate counters, together with the standard packages those lines need.  The retained environments print in this document as Theorem 1 in order of appearance, whereas the article prints the same items with the numbering of its own full document.  The complete native source of the article as distributed in the arXiv e-print for this exact version is bundled unchanged as \texttt{sources/191728/source0.tex}.
\begin{theorem}\label{thm: associativity}
    The $K$-Hall algebra $\BK(\CA)$ is an associative algebra.
\end{theorem}

\begin{proof}
We fix
\[\varphi_\alpha\in K_\ast(\FM^\rig_\alpha)\,,\,\varphi_\beta\in K_\ast(\FM^\rig_\beta)\,,\,\varphi_\gamma\in K_\ast(\FM^\rig_\gamma)\,.\]
We will use $\alpha, \beta, \gamma$ as indices in the different objects involved in the definition of $\ast$ to keep track of which copy we refer to. For example, $\Sigma_{\alpha, \beta}$ refers to the restriction of the direct sum map to $\FM_\alpha\times \FM_\beta\to \FM_{\alpha+\beta}$ and $\Gamma_{\alpha, \beta}$ refers to the restriction of $\Gamma_-$ to $\FM_\alpha\times \FM_\beta$ or $(\FM_\alpha\times \FM_\beta)^\rig$; we ease the notation by writing only $\Sigma_{\alpha, \beta}$ instead of $\Sigma^\rig_{\alpha, \beta}$. To calculate $\varphi_\alpha\ast(\varphi_\beta\ast \varphi_\gamma)$ we consider the following diagram, where the square is cartesian:
\begin{center}
\begin{tikzcd}
    (\FM_\alpha\times \FM_\beta\times \FM_\gamma)^\rig\arrow[r, "\overline \pi_{\alpha, \beta+\gamma}"]  \arrow[dd, bend right=70, swap,"\Sigma_{\alpha, \beta, \gamma}"]\arrow[rr, bend left=20, "\pi_{\alpha, \beta, \gamma}"]\arrow[d, "\overline \Sigma_{\beta, \gamma}"]&
    \FM_\alpha^\rig\times (\FM_\beta\times \FM_\gamma)^\rig \arrow[r, "\pi_{\beta,\gamma}"]\arrow[d, "\Sigma_{\beta, \gamma}"]&
    \FM_\alpha^\rig\times \FM_\beta^\rig\times \FM_\gamma^\rig\\
     (\FM_\alpha\times \FM_{\beta+\gamma})^\rig \arrow[r, "\pi_{\alpha, \beta+\gamma}"] \arrow[d, "\Sigma_{\alpha, \beta+\gamma}"]&
     \FM_\alpha^\rig\times \FM_{\beta+\gamma}^\rig & \\
     \FM_{\alpha+\beta+\gamma}^\rig
\end{tikzcd}
\end{center}
The map $\Sigma_{\alpha, \beta, \gamma}$ is the triple direct sum map $\FM_{\alpha}\times \FM_\beta\times \FM_\gamma\to \FM_{\alpha+ \beta+ \gamma}$ (or its rigidification), and by associativity of the direct sum it can be written as
\[\Sigma_{\alpha, \beta, \gamma}\coloneqq \Sigma_{\alpha, \beta+\gamma}\circ \Sigma_{\beta,\gamma}=\Sigma_{\alpha+\beta, \gamma}\circ\Sigma_{\alpha, \beta}\,.\]
The map $\pi_{\alpha, \beta, \gamma}$ is the rigidification with respect to the natural $BT$ acting on  $(\FM_\alpha\times \FM_\beta\times \FM_\gamma)^\rig$ -- more precisely, $T$ is the 2-dimensional torus obtained by modding out the diagonal $\BG_m$ from $\BG_m^3$. We have the associativity type relation
\[\pi_{\alpha, \beta, \gamma}\coloneqq \pi_{\beta, \gamma}\circ \overline \pi_{\alpha, \beta+\gamma}=\pi_{\alpha, \beta}\circ \overline \pi_{\alpha+\beta, \gamma}\,.\]

Denote $\varphi_{\alpha\beta\gamma}=\varphi_\alpha\boxtimes\varphi_\beta\boxtimes \varphi_\gamma\in K_\ast(\FM_\alpha^\rig\times \FM_\beta^\rig\times \FM_\gamma^\rig)$. Then
\begin{align}\label{eq: proofassociativity}
    \varphi_\alpha\ast(\varphi_\beta\ast \varphi_\gamma)&=(\Sigma_{\alpha, \beta+\gamma})_\ast\left(\Gamma_{\alpha, \beta+\gamma}\otimes \pi_{\alpha, \beta+\gamma}^\ast(\Sigma_{\beta, \gamma})_\ast\left(\Gamma_{\beta, \gamma}\otimes \pi_{\beta, \gamma}^\ast(\varphi_{\alpha\beta\gamma})\right)\right)\\
    &\nonumber =(\Sigma_{\alpha, \beta+\gamma})_\ast\left(\Gamma_{\alpha, \beta+\gamma}\otimes (\overline\Sigma_{\beta, \gamma})_\ast\overline\pi_{\alpha, \beta+\gamma}^\ast\left(\Gamma_{\beta, \gamma}\otimes \pi_{\beta, \gamma}^\ast(\varphi_{\alpha\beta\gamma})\right)\right)\\
    &\nonumber =(\Sigma_{\alpha, \beta, \gamma})_\ast \left(\overline \Sigma_{\beta, \gamma}^\ast(\Gamma_{\alpha, \beta+\gamma})\otimes \overline\pi_{\alpha, \beta+\gamma}^\ast(\Gamma_{\beta, \gamma})\otimes \pi_{\alpha, \beta, \gamma}^\ast(\varphi_{\alpha\beta\gamma})\right)
\end{align}

Note that, by the properties of the Ext complex, we have
\begin{align*}\Ext_{<}&\coloneqq \Ext_{\alpha, \beta}+\Ext_{\alpha, \gamma}+\Ext_{\beta, \gamma}=\Sigma_{\beta,\gamma}^\ast(\Ext_{\alpha, \beta+\gamma})+\Ext_{\beta, \gamma}\\
\Ext_{>}&\coloneqq \Ext_{\beta, \alpha}+\Ext_{\gamma, \alpha}+\Ext_{\gamma, \beta}=\Sigma_{\beta,\gamma}^\ast(\Ext_{\beta+\gamma, \alpha})+\Ext_{ \gamma, \beta}
\end{align*}
on $\FM_{\alpha}\times \FM_\beta\times \FM_\gamma$, and therefore
\[\overline \Sigma_{\beta, \gamma}^\ast(\Gamma_{\alpha, \beta+\gamma})\otimes \overline\pi_{\alpha, \beta+\gamma}^\ast(\Gamma_{\beta, \gamma})=\Lambda_{-1}\big(\Ext_{<}+\Ext_{>}^\vee\big)\otimes \det(\Ext_{<})^\vee[\rk \Ext_{<}]\,,\]
which we can plug in \eqref{eq: proofassociativity}. An entirely analogous calculation arrives at the same expression for $(\varphi_\alpha\ast \varphi_\beta)\ast \varphi_\gamma$. 
\end{proof}

\end{document}
